\chapter{Conclusion and Future Work}
\label{ch:conclusion}

\section{Conclusion}

This project set out to answer a narrow, specific question: can a PPO
agent, trained in real time against a live Kubernetes cluster, match or
improve on a correctly tuned Horizontal Pod Autoscaler on held-out traffic,
under a decision rule fixed before any evaluation result exists? Across two
generations of design, the honest answer is that we cannot declare a
winner, and we think that answer is itself a useful result rather than a
disappointing one, for three reasons.

First, V1's failure was informative, not just negative. One of three
required valid PPO repeats completing is not "PPO lost a close contest";
it is a sign that the controller's own reward function and replica floor
made it structurally fragile in a way the evaluation correctly exposed. We
traced that fragility to two specific, nameable design decisions -- an
unconditional per-replica resource cost and a replica floor of one -- and
built V2 specifically to remove them, rather than retraining V1's design
for longer and hoping the problem would go away on its own.

Second, V2's pre-registered evaluation is a genuinely negative result for
the strong form of the hypothesis ("RL beats HPA"), not because PPO
performed badly, but because the evaluation's own traffic load was light
enough that the primary latency objective sat at or near its floor for
every arm, and the metric that did discriminate between arms, cost, did not
favour any single arm across all three profiles. We regard this as the
evaluation doing its job rather than as an inconclusive accident: the
pre-registered decision rule was built specifically to stop us from
declaring a winner on weak or profile-dependent evidence, and it did
exactly that, on our own result.

Third, the exploratory stress test, run at training-level peak loads where
the latency objective does separate arms, still found no controller that
was uniformly better across all three stress profiles -- PPO was cheapest
and had the lowest latency under a CPU burst, but had the highest
SLA-violation rate under a mixed burst, the opposite ordering from the CPU
case. This is consistent with a picture in which each controller's
relative strength depends on which resource dimension (CPU or memory) is
actually under pressure, which is a more specific and more useful finding
than either "RL wins" or "RL loses" would have been.

We also want to state plainly what this report does \emph{not} claim. It
does not claim that RL-based autoscaling cannot outperform HPA in general;
our evaluation loads were light relative to training-level peaks, our
sample sizes were small (three repeats for the pre-registered campaign, one
repeat per cell for the exploratory stress test), and our testbed was a
single-node cluster on a laptop. It does not claim that the V2 controller
is production-ready; no safety termination occurred in either the
pre-registered evaluation or the stress test, but that absence of failure
over 36 total runs is not the same as a guarantee over a much longer
operational horizon. What it does claim is that we built and ran a
complete, live, pre-registered comparison end to end, including the
failures along the way, and reported what actually happened rather than
what we had hoped would happen.

\section{Future Work}

Several directions follow directly from the limitations in
Section~\ref{sec:limitations}. A pre-registered evaluation at
training-level peak loads, with the same three-repeats-per-cell discipline
used for \texttt{eval\_\allowbreak final} rather than the single-repeat, descriptive
design of \texttt{stress\_\allowbreak explore}, would let the decision rule actually be
applied at the load levels where the primary metric is not stuck at its
floor. A longer training budget, beyond the 12,288 timesteps used here,
would let us check whether explained variance and entropy, which were still
trending in the expected direction at the end of training rather than
having clearly plateaued, continue to improve with more experience. Running
the same comparison on a multi-node cluster, where node-level resource
contention and pod scheduling behave differently from a single-node
Minikube cluster, would test whether our findings transfer beyond the
testbed used here. Finally, the histogram-bucket artefact at 975~ms that we
flag throughout Chapter~\ref{ch:results} points to a concrete, fixable
instrumentation improvement -- a finer-grained latency histogram, or a
summary-based latency estimator -- that would let future work report p95
latency on these memory- and mixed-mode profiles with less ambiguity than
our own figures currently can.
